\documentclass[10pt,a4paper]{amsart}
\usepackage[margin=19mm]{geometry}
\usepackage{amsmath,amssymb,mathtools}
\usepackage[scr=rsfs,cal=euler]{mathalfa}
\usepackage{xfrac}
\usepackage[unicode,hidelinks,bookmarks=false]{hyperref}
\newcommand{\ie}{i.e.\ }
\newcommand{\cf}{cf.\ }
\newcommand{\resp}{respectively\ }
\newcommand{\Subsection}{\subsection}
\providecommand{\nobreakdash}{\nobreak\hbox{-}}
\setlength{\emergencystretch}{2em}
\multlinegap0pt
\usepackage{bm}
\usepackage{bbm}
\usepackage{dsfont}
\usepackage{xspace}
\usepackage{cancel}
\usepackage{mathtools}
\usepackage{amsthm}
\makeatletter
\@ifundefined{lemma}{\newtheorem{lemma}{Lemma}}{}
\makeatother
\DeclareMathOperator{\Det}{Det}
\DeclareMathOperator{\adj}{adj}
\DeclareMathOperator{\rank}{rank}
\newcommand{\CC}{\mathbb{C}}
\newcommand{\Mcal}{\mathcal M}
\newcommand{\OO}{\mathcal O}
\newcommand{\PP}{\mathbb{P}}
\title[A symmetric determinantal lower bound for diagonal power sum]{A symmetric determinantal lower bound for diagonal power sums via polar degree}
\author{Karthik Sheshadri}
\date{Source: arXiv:2606.11090v1 (CC BY 4.0)}
\begin{document}
\maketitle

\noindent\emph{Source and licence.} A symmetric determinantal lower bound for diagonal power sums via polar degree. Version arXiv:2606.11090v1 (CC BY 4.0), distributed under the Creative Commons Attribution 4.0 International licence, as recorded in the arXiv licence metadata for this exact version and shown on the abstract page https://arxiv.org/abs/2606.11090v1. Full rights notice: licenses/arxiv-2606.11090-CC-BY-4.0.txt.\par\smallskip

\noindent\textit{Editorial scope.} Counted unit: the Lemma whose source label is \texttt{lem:isolated} in the article (source0.tex lines 396--410, source label \texttt{lem:isolated}), delivered together with the article's own complete original proof of it (source0.tex lines 412--511, one \texttt{proof} environment of 100 lines). No mathematics is written, repaired or re-derived: statement and proof are the article's own text line for line. Required context retained uncounted: eq:homdet (lines 301--304); lem:rank (lines 317--325). Every definition, display and label of those lines is retained unchanged. Omitted from this package and not redistributed: 1--300, 305--316, 326--395, 411--411, 512--989. Nothing inside a retained excerpt is omitted or altered: every retained body is delivered exactly as the article's lines are written, so no omission receipt of any kind accompanies this package. Citations: the retained excerpts carry no citation command at all, so no bibliography entry of the article is transcribed and no reference is redistributed. Reference adaptation: none was needed and none was made; every reference inside the delivered unit resolves inside it, so no unresolved reference or citation is produced. Printed slips kept literal: no printed word, symbol, hypothesis or number of the counted statement or of its proof was altered. Layout and class: the article's own class is \texttt{article} with packages including geometry, amsmath, amssymb, amsthm, mathtools, bm, hyperref, enumitem, microtype, upquote; the wrapper is the lane's pinned \texttt{amsart} editorial wrapper at 10pt on A4, which restates the theorem-like environments the retained text needs and adds the article's own macro definitions that the retained text needs (\texttt{\textbackslash CC}, \texttt{\textbackslash Det}, \texttt{\textbackslash Mcal}, \texttt{\textbackslash OO}, \texttt{\textbackslash PP}, \texttt{\textbackslash adj}, \texttt{\textbackslash rank}), with extra wrapper packages (bm, bbm, dsfont, xspace, cancel, mathtools). The delivered numbering is the unit's own. Assets: no retained span contains a figure environment, a graphics-inclusion command, a PDF-page inclusion, a table or a TikZ picture; no image file is copied into this package, the package declares no asset dependency and no image file is redistributed with it. Licence: version arXiv:2606.11090v1 of this article is distributed under the Creative Commons Attribution 4.0 International licence, as recorded in the arXiv licence metadata for this exact version and shown on the abstract page https://arxiv.org/abs/2606.11090v1. A full rights notice is delivered at licenses/arxiv-2606.11090-CC-BY-4.0.txt. Characters: every retained excerpt is pure ASCII, so the delivered engine, XeLaTeX with its default Latin Modern text font, reports no missing character.
\begin{equation}
\label{eq:homdet}
  \Det\Mcal(z,x)=z^{m-d}f(x).
\end{equation}

\begin{lemma}[Rank at genuine smooth points]
\label{lem:rank}
Let $[z_0:x_0]\in\PP^N$ satisfy $z_0\ne0$ and $f(x_0)=0$, with
$[x_0]\in X$ smooth.  Then
\[
  \rank\Mcal(z_0,x_0)=m-1.
\]
Consequently the projective kernel line $[u_0]\in\PP^{m-1}$ is unique.
\end{lemma}

\begin{lemma}[Symmetric isolatedness]
\label{lem:isolated}
Let $p=([z_0:x_0],[u_0])$ be the lift of a point of $P_q$ to the symmetric
incidence scheme
\begin{equation}
\label{eq:incidence}
  \Mcal(z,x)u=0,
  \qquad
  q_j(u^{\mathsf T}A_1u,\ldots,u^{\mathsf T}A_Nu)=0\ (1\le j\le N-2),
  \qquad
  h(z,x)=0
\end{equation}
in $\PP^N_{[z:x]}\times\PP^{m-1}_{[u]}$.  Then $p$ is a zero-dimensional
scheme-theoretically isolated solution of \eqref{eq:incidence}.
\end{lemma}

\begin{proof}
By Lemma~\ref{lem:rank}, $\Mcal(z_0,x_0)$ has rank $m-1$.  Since it is
symmetric, its kernel is a one-dimensional radical.  After a constant congruence
change of basis $\Mcal\mapsto P^{\mathsf T}\Mcal P$, which gives an
isomorphic local incidence problem via $u=P u'$ and preserves the quadratic
forms since ${u'}^{\mathsf T}(P^{\mathsf T}A_iP)u'=u^{\mathsf T}A_i u$, we
may assume that the kernel line at $p$ is spanned by the last coordinate vector.
The determinant is multiplied only by the nonzero scalar $(\det P)^2$, which is
harmless in all local unit comparisons below.  Because the kernel is the radical
of the symmetric bilinear form, the restriction to any complementary
$(m-1)$-dimensional subspace is nondegenerate.  Thus, in a Zariski neighborhood
of $[z_0:x_0]$, write
\[
  \Mcal=\begin{pmatrix}B&c\\ c^{\mathsf T}&s\end{pmatrix},
  \qquad \det B\in\OO^\times.
\]
Work on the affine chart $u_m=1$ and write $u=(u',1)^{\mathsf T}$.  The kernel
equations are
\[
  Bu'+c=0,
  \qquad
  c^{\mathsf T}u'+s=0.
\]
Since $B$ is invertible in the local ring, the first $m-1$ equations are
equivalent to
\[
  u'=-B^{-1}c.
\]
With
\[
  g=s-c^{\mathsf T}B^{-1}c,
\]
the last equation becomes
\[
  c^{\mathsf T}u'+s
  =c^{\mathsf T}(u'+B^{-1}c)+g.
\]
Therefore the ideal generated by $\Mcal u$ is exactly
\[
  \big(u'+B^{-1}c,\ g\big)
\]
in the local ring.  Scheme-theoretically, the kernel incidence is the graph
\[
  u=\nu:=\begin{pmatrix}-B^{-1}c\\1\end{pmatrix}
\]
over the hypersurface $g=0$.  The determinant is
\begin{equation}
\label{eq:schurdet}
  \Det\Mcal=(\det B)g.
\end{equation}

On $g=0$, the matrix $\Mcal$ has kernel spanned by $\nu$ and rank $m-1$ because
$B$ remains invertible.  Its adjugate is therefore a scalar multiple of
$\nu\nu^{\mathsf T}$.  The lower-right cofactor equals $\det B$, while the
lower-right entry of $\nu\nu^{\mathsf T}$ is $1$, so
\begin{equation}
\label{eq:adjlocal}
  \adj(\Mcal)=(\det B)\nu\nu^{\mathsf T}
  \qquad\text{on }g=0.
\end{equation}
Hence, in the local quotient by the kernel equations,
\[
  \frac{\partial}{\partial x_i}\Det\Mcal
  =(\det B)\nu^{\mathsf T}A_i\nu.
\]
Using \eqref{eq:homdet}, we get
\begin{equation}
\label{eq:unitpolar}
  \nu^{\mathsf T}A_i\nu
  =\frac{z^{m-d}}{\det B}\,\partial_i f(x)
  \qquad\text{on the local kernel graph.}
\end{equation}
The factor $z^{m-d}/\det B$ is a unit near $p$, since $z_0\ne0$ and
$\det B$ is a unit.

It follows that, after eliminating the kernel variables, each lifted polar
equation
\[
  q_j(\nu^{\mathsf T}A_1\nu,\ldots,\nu^{\mathsf T}A_N\nu)=0
\]
is the ordinary polar equation
\[
  q_j(\partial_1f,\ldots,\partial_Nf)=0
\]
multiplied by the same local unit.  By \eqref{eq:schurdet} and
\eqref{eq:homdet}, the equation $g=0$ is also the equation $f=0$ up to a unit on
$z\ne0$.

Thus the completed local ring of the incidence scheme at $p$ is isomorphic to
\[
  \widehat\OO_{\PP^N,[z_0:x_0]}
  \Big/
  \big(f,\ q_1(\nabla f),\ldots,q_{N-2}(\nabla f),\ h\big),
\]
up to multiplication of the displayed equations by units.  The chosen
$q_j$ cut out the finite reduced polar set $P_q$ on $X$, and the slice $h$ meets
the corresponding cone line at exactly one point.  Hence this local ring is
zero-dimensional.  Since the schemes are of finite type over $\CC$, it is
Artinian and has finite length.
\end{proof}

\end{document}
